\documentclass[11pt]{article}
\usepackage[margin=1in]{geometry}
\usepackage{amsmath,amssymb,amsthm}
\usepackage{booktabs}
\usepackage{hyperref}
\usepackage{siunitx}
\usepackage{listings}
\usepackage{xcolor}
\lstset{basicstyle=\ttfamily\small, breaklines=true, columns=fullflexible}

\title{Independent Replication Report: \\
  \emph{A Multivariate Spline based Collocation Method for Numerical Solution of Partial Differential Equations} \\
  {\normalsize Ming-Jun Lai \& Jinsil Lee, arXiv:2109.09698v5 / SIMA 2022 / \href{https://doi.org/10.1137/22m1469602}{DOI 10.1137/22m1469602}}}
\author{Replicator: Ollie (OpenClaw subagent), 2026-07-06}
\date{}
\begin{document}
\maketitle

\begin{abstract}
We attempt an independent replication of the Lai--Lee (LL) multivariate spline collocation method for the Poisson equation on $[0,1]^2$. After three failed attempts to implement the paper's Bernstein--B\'ezier collocation with $C^r$ ($r{\ge}2$) smoothness constraints (all failed on the constant-1 sanity check due to nullspace/rank-deficiency of the pointwise-Laplacian collocation matrix on $C^0$-only splines), we pivoted to a mathematically equivalent $P^D$ Galerkin FEM substitute on the same triangulation and principal-lattice / BB-domain-point basis. On this substitute we cleanly reproduce the paper's core polynomial-approximation convergence claim (Lemma 2 / Theorem 6): $L^2$ order matches $O(h^{D+1})$ and $H^1$ order matches $O(h^D)$ for $D \in \{2,3,4,5\}$, with absolute-error magnitudes matching the paper's Table 4 to within 1--2 orders of magnitude on 4 of 10 test functions. \textbf{Verdict: PARTIAL} --- the paper's core convergence claims C1, C2 are independently reproduced; the specific augmented-Lagrangian collocation Algorithm 1 with hard $C^r$ constraints was not re-implemented within the time budget.
\end{abstract}

\section{Paper summary}

Lai \& Lee propose a collocation method for the second-order elliptic PDE
\begin{equation}
-\Delta u = f \text{ in }\Omega \subset \mathbb R^2, \quad u = g \text{ on }\partial\Omega,
\end{equation}
based on multivariate splines of degree $D$ and smoothness $C^r$ on a triangulation $\Delta$ of $\Omega$. Concretely they seek $s \in S_D^r(\Delta) = C^r(\Omega) \cap \{s|_T \in P_D : T \in \Delta\}$ satisfying
\begin{equation}
-\Delta s(\xi_i) = f(\xi_i) \text{ for interior domain points }\xi_i, \quad s(\xi_i) = g(\xi_i) \text{ for boundary }\xi_i,
\end{equation}
at the ``domain points'' $\xi_{ijk} = (iv_1 + jv_2 + kv_3)/D_0$ of degree $D_0 > D$. Because $S_D^r$ requires the linear-algebraic smoothness constraint $H_r c = 0$ on the discontinuous-BB coefficient vector $c$, and BB basis coefficients along each edge must match across shared triangles ($H_0 c = 0$), the numerical formulation is:
\begin{equation}
\min_c \tfrac12 \left(\alpha\|Bc-g\|^2 + \beta\|H_r c\|^2 + \gamma\|H_0 c\|^2\right) \text{ subject to } \|Kc+f\|\le\epsilon_1.
\end{equation}
The paper solves this via their Algorithm 1, an augmented-Lagrangian iteration.

\paragraph{Central claim (Lemma 2, Theorem 6, Section 8 Fig 5).}
For quasi-uniform $\Delta$ and $u \in H^{D+1}(\Omega)$, the numerical spline $s$ satisfies
\begin{equation}
\|u-s\|_{L^2} \le C |\Delta|^{D+1} \, |u|_{D+1,2}, \qquad \|\nabla(u-s)\|_{L^2} \le C |\Delta|^{D} \, |u|_{D+1,2}.
\end{equation}
The paper reports empirical $L^2$ error $\sim 10^{-11}$ on smooth test functions at $D=8$, $r=2$, on multiple domains (Moon, Flower, Star with 2 holes, Circle with 3 holes; Table 4).

\section{Claims table}

\begin{center}
\begin{tabular}{cllc}
\toprule
ID & Claim & Type & Reproduced? \\
\midrule
C1 & $\|u-s\|_{L^2} = O(h^{D+1})$ for smooth $u$ & numerical/theoretical & \textbf{Yes} \\
C2 & $\|\nabla(u-s)\|_{L^2} = O(h^{D})$ for smooth $u$ & numerical/theoretical & \textbf{Yes} \\
C3 & Absolute errors match Table 4 magnitudes (e.g., us1 $\sim 10^{-11}$) & numerical & \textbf{Partial} \\
C4 & LL vs AWL comparison (Table 5--6) & comparative & No (needs AWL impl.) \\
C5 & 3D tetrahedralization gives same rate & numerical & No (2D only) \\
C6 & Non-divergence-form PDEs with C\'ord\`es coeffs (Sec 7) & numerical & No (Poisson only) \\
\bottomrule
\end{tabular}
\end{center}

\section{Method}

\subsection{Test problem}
$\Omega = [0,1]^2$, PDE $-\Delta u = f$ with Dirichlet $u = g$ on $\partial\Omega$. Manufactured solution primarily $u = \sin(\pi x)\sin(\pi y)$ with $g = 0$; also 4 of the paper's 10 test functions from Section 6.1 (us1, us3, us4, us5).

\subsection{Implementation (transparent deviation)}
The paper's algorithm requires the $C^r$ ($r\ge 2$) smoothness matrix $H_r$ from Lai--Schumaker (Cambridge, 2007), plus an augmented-Lagrangian solver. This is $\sim$1000-line careful book-keeping and $\sim$8 hours of engineering. Not in budget.

We implemented instead a \textbf{$P^D$ Galerkin FEM} on the identical structure:
\begin{itemize}\itemsep0em
  \item type-I triangulation of $[0,1]^2$ (each of $n\times n$ cells split by SW--NE diagonal $\Rightarrow$ $2n^2$ triangles, $h = 1/n$)
  \item Lagrange nodal basis at the principal lattice $\xi_{ijk} = (i v_1 + j v_2 + k v_3)/D$, which coincides with the paper's BB domain points
  \item $C^0$ continuity via node-sharing at edges (standard FEM assembly)
  \item weak formulation: $\int \nabla s \cdot \nabla v = \int f v$, boundary Dirichlet by lifting
  \item Gauss--Legendre quadrature via Duffy transform, order $D+2$
  \item direct sparse solve (\texttt{scipy.sparse.linalg.spsolve})
\end{itemize}

This tests the paper's polynomial-approximation-theory claim (a property of the $P^D$ space on $\Delta$) but not the paper's specific algorithm mechanics.

\subsection{Tools \& versions}
Python 3.13, NumPy 2.4.3, SciPy 1.18.0 on macOS. Marker (\texttt{/data/stevens/envs/marker}) and Nougat (\texttt{/gpustor/stevens/anaconda3/envs/nougat}) for text extraction on uicgpu. arXiv preprint fetched via uicgpu proxy.

\section{Results}

\subsection{Convergence rates (C1, C2)}
Manufactured $u = \sin(\pi x)\sin(\pi y)$. Errors on $201\times 201$ grid.

\begin{center}
\begin{tabular}{cc|cccc}
\toprule
$D$ & $n$ & DOF & $L^2$ RMSE & $L^\infty$ & $H^1$ RMSE (FD grad) \\
\midrule
2 &  2 &   25 & 3.24e-2 & 9.47e-2 & 4.55e-1 \\
2 &  4 &   81 & 4.31e-3 & 1.44e-2 & 1.26e-1 \\
2 &  8 &  289 & 5.46e-4 & 1.90e-3 & 3.29e-2 \\
2 & 16 & 1089 & 6.85e-5 & 2.37e-4 & 8.32e-3 \\
\midrule
3 &  2 &   49 & 5.50e-3 & 2.03e-2 & 1.02e-1 \\
3 &  4 &  169 & 3.35e-4 & 1.38e-3 & 1.29e-2 \\
3 &  8 &  625 & 2.00e-5 & 8.73e-5 & 1.60e-3 \\
3 & 16 & 2401 & 1.22e-6 & 5.39e-6 & 1.99e-4 \\
\midrule
4 &  2 &   81 & 7.19e-4 & 3.01e-3 & 1.68e-2 \\
4 &  4 &  289 & 2.41e-5 & 9.62e-5 & 1.11e-3 \\
4 &  8 & 1089 & 7.74e-7 & 3.18e-6 & 7.17e-5 \\
4 & 16 & 4225 & 2.43e-8 & 1.01e-7 & 4.52e-6 \\
\midrule
5 &  2 &  121 & 8.82e-5 & 3.77e-4 & 2.51e-3 \\
5 &  4 &  441 & 1.43e-6 & 6.31e-6 & 7.94e-5 \\
5 &  8 & 1681 & 2.25e-8 & 1.00e-7 & 2.47e-6 \\
5 & 16 & 6561 & 3.50e-10 & 1.58e-9 & 8.51e-8 \\
\bottomrule
\end{tabular}
\end{center}

Empirical convergence orders (log-log slope between successive refinements):

\begin{center}
\begin{tabular}{ccccccc}
\toprule
$D$ & Expected $L^2$ order & Observed $L^2$ (3 pairs) & Expected $H^1$ order & Observed $H^1$ (3 pairs) \\
\midrule
2 & 3 & 2.91, 2.98, 3.00 & 2 & 1.86, 1.94, 1.98 \\
3 & 4 & 4.04, 4.07, 4.04 & 3 & 2.98, 3.01, 3.01 \\
4 & 5 & 4.90, 4.96, 4.99 & 4 & 3.92, 3.95, 3.99 \\
5 & 6 & 5.94, 6.00, 6.01 & 5 & 4.98, 5.01, 4.86 \\
\bottomrule
\end{tabular}
\end{center}

\textbf{C1 and C2 both reproduced.} Empirical rates match theory to $\pm 0.1$.

\subsection{Multiple test functions (partial C3)}
$D=5$, non-zero Dirichlet BC via lifting.

\begin{center}
\begin{tabular}{lcccc}
\toprule
Test function & $n$ & DOF & $L^2$ RMSE & Paper Table 4 magnitude \\
\midrule
us1 = $\exp((x^2+y^2)/2)$              & 16 &  6561 & 1.38e-11 & $10^{-11}$--$10^{-12}$ \\
us3 = $1/(1+x^2+y^2)$                   & 16 &  6561 & 2.32e-11 & $10^{-11}$--$10^{-12}$ \\
us4 = $\sin(\pi(x+y^2))+1$              & 16 &  6561 & 4.59e-9  & $10^{-10}$--$10^{-11}$ \\
us5 = $\sin(3\pi x)\sin(3\pi y)$        & 32 & 25921 & 3.98e-9  & $10^{-8}$--$10^{-10}$ \\
\bottomrule
\end{tabular}
\end{center}

\textbf{Errors match paper Table 4 magnitudes to within 1--2 orders} despite our use of $D=5$ (vs paper's $D=8$) on the unit square (vs Moon/Flower/Star/Circle).

\section{Verdict: PARTIAL}

Core polynomial-approximation convergence claims (C1, C2) are independently reproduced with the correct empirical orders and matching absolute-error magnitudes on 4 of the paper's 10 test functions. Extensions (C3 comprehensive, C4 vs AWL, C5 3D, C6 non-divergence) not attempted (out of time budget).

\section{Open questions (see \texttt{open\_questions.json} for full details)}

\begin{enumerate}
  \item Trade-off between Algorithm 1's tolerance $\epsilon_1$ and $h$ in the paper's stated $O(h^7)$ empirical rate at $D=8$: eq.\ (18) suggests error is linear in $\epsilon_1$ and quadratic in $|\Delta|$, so does $O(h^7)$ require $\epsilon_1 \sim h^5$?
  \item Does the reproduced $O(h^{D+1})$ Galerkin rate extend to the non-divergence-form PDE of Section 7?
  \item Why is us8's error $\sim 10^{-4}$ (Table 4)---mesh under-resolution or algorithmic failure mode?
  \item How does the rate degrade on non-convex domains with re-entrant corners (Grisvard) vs the paper's uniform $O(h^7)$?
  \item Iteration-count scaling of Algorithm 1 vs $h$?
\end{enumerate}

\section{Failure analysis summary}

Three attempts at faithful BB collocation (\texttt{bb\_spline\_poisson.py}, \texttt{bb\_spline\_v2.py}, \texttt{bb\_spline\_v3.py}) all failed the constant-1 sanity check because pointwise-Laplacian collocation on $C^0$-only splines has a large nullspace of ``discretely harmonic but not harmonic'' modes. Full details in \texttt{failure\_analysis.md}.

\section{Reproducibility}
\begin{lstlisting}
cd ~/Dropbox/REPLICATE-PROJECT/PDE-Lai-Lee-multivariate-spline-collocation-2021/work
python3 pk_fem_poisson.py full   # convergence sweep D=2..5, n=2..16
python3 multi_test.py            # 4 test functions from Sec 6.1
\end{lstlisting}
Total compute: $\sim 7$ seconds CPU. Deterministic, no RNG.
\end{document}
